%! AP Statistics — Unit 4, Lesson 4.11 — Group Activity KEY
\documentclass[10pt]{article}
\usepackage{apstats-article}
\usepackage{apstats-key}

\begin{document}

\pageheader{Unit 4, Lesson 4.11}{Group Activity KEY: Binomial Parameters in Context}
\namepartnerperiod

\vspace{0.1in}

% ── TIER R ────────────────────────────────────────────────────────────────────
\begin{tcolorbox}[colback=goldbg, colframe=goldacc, boxrule=1pt, arc=2mm,
    title={\bfseries Tier R --- Key},
    fonttitle=\small]
\small

\textbf{(a)} $X \sim B(25, 0.40)$

$\mu_X = 25(0.40) = \ans{10}$

$\sigma_X = \sqrt{25(0.40)(0.60)} = \sqrt{6} \approx \ans{2.45}$

\ans{If this binomial experiment were repeated many times, the average number of successes would be 10.}

\vspace{10pt}
\textbf{(b)} $X \sim B(100, 0.08)$

$\mu_X = 100(0.08) = \ans{8}$

$\sigma_X = \sqrt{100(0.08)(0.92)} = \sqrt{7.36} \approx \ans{2.71}$

\ans{In many repetitions of this 100-trial experiment, the average number of successes would be 8.}
\end{tcolorbox}

\vspace{0.1in}

% ── TIER A ────────────────────────────────────────────────────────────────────
\begin{tcolorbox}[colback=sky, colframe=navy, boxrule=1pt, arc=2mm,
    title={\bfseries Tier A --- Key},
    fonttitle=\small]
\small

Call center: $n=80$, $p=0.22$.

(a) $\mu_X = 80(0.22) = \ans{17.6}$; \quad $\sigma_X = \sqrt{80(0.22)(0.78)} = \sqrt{13.728} \approx \ans{3.71}$

(b) \ans{If the call center handles many afternoons of 80 calls each, it would average 17.6 first-attempt resolutions per afternoon.}

(c) $P(X \le 15) = \texttt{binomcdf(80, 0.22, 15)} \approx \ans{0.2752}$

(d) $P(X \ge 20) = 1 - P(X \le 19) = 1 - \texttt{binomcdf(80, 0.22, 19)} \approx \ans{1 - 0.7305 = 0.2695}$

(e) \ans{$\mu_X + 2\sigma_X = 17.6 + 2(3.71) = 25.02$. Since 20 is within two standard deviations of the mean, it is not unusually high. It is about 0.65 standard deviations above the mean.}
\end{tcolorbox}

\vspace{0.1in}

% ── TIER E ────────────────────────────────────────────────────────────────────
\begin{tcolorbox}[colback=greenbg, colframe=greenacc, boxrule=1pt, arc=2mm,
    title={\bfseries Tier E --- Key},
    fonttitle=\small]
\small

(a)
\renewcommand{\arraystretch}{1.8}
\begin{tabularx}{\linewidth}{@{} >{\bfseries}p{2cm} X X @{}}
 & Process A & Process B \\
\hline
$\mu_X$ & \ans{$200(0.05)=10$} & \ans{$100(0.10)=10$} \\
$\sigma_X$ & \ans{$\sqrt{200(0.05)(0.95)} \approx 3.08$} & \ans{$\sqrt{100(0.10)(0.90)} = 3.00$} \\
\end{tabularx}

(b) \ans{Both have $\mu_X = 10$ because $np$ is the same: $200 \times 0.05 = 100 \times 0.10 = 10$. A lower defect rate with more trials can produce the same expected count as a higher defect rate with fewer trials.}

(c) \ans{Process A has slightly more variability ($\sigma_X \approx 3.08$ vs.\ $3.00$). Practically, the counts of defects in Process A will fluctuate slightly more from batch to batch, making quality control marginally harder to predict.}

(d) \ans{$\sigma_X = \sqrt{np(1-p)}$ is maximized when $p(1-p)$ is maximized. Taking the derivative: $\frac{d}{dp}[p(1-p)] = 1-2p = 0 \Rightarrow p = 0.5$. At $p=0.5$, $p(1-p) = 0.25$ is its maximum, so $\sigma_X$ is maximized at $p = 0.5$.}
\end{tcolorbox}

\end{document}
